\section{Bounded specialization and stabilization of the minimum}
\label{sec:bounded}

\subsection{Arithmetic setup and spreading the class}
Let $K\subset\C$ be a number field, let $U$ be a nonempty open subscheme of
$\Spec\mathcal O_K$, and let $f:X\to U$ be smooth and projective of pure relative
dimension $g$. Fix a relatively ample line bundle $L$, an integer $0\leq r\leq g$,
and put $d=g-r$. All special-fiber cycles below are geometric cycles. For an
effective $d$-cycle $Z$, with integer multiplicities, write
\[
 \deg_L Z=c_1(L)^d\cdot Z.
\]
If a projective embedding is needed, choose one relatively very ample power
$L^m$. The embedding degree is then $m^d\deg_L Z$; the complexity throughout
continues to use the original $L$-degree.

Let $t$ be a rational absolute Hodge class on the generic fiber. Deligne,
\emph{Hodge cycles on abelian varieties}, Proposition~2.9(b), says that the
Galois action on the space of absolute Hodge classes factors through a finite
quotient. After a finite extension of $K$, we may therefore assume that $t$ is
fixed. Choose one auxiliary prime $\ell$ and remove the primes above $\ell$
from $U$. Smooth proper base change makes
\[
 \mathcal V_\ell=R^{2r}f_*\Q_\ell(r)
\]
a lisse sheaf. Its generic Galois-invariant vector $t_\ell$ gives a global
section: the representation is unramified on $U$, and the generic Galois group
surjects onto the relevant fundamental group. Finite extension of $K$ does not
change the geometric minima, and inverse images of a finite set of finite
places are finite. Thus the conclusions below descend to the original
number field.

For a geometric point $s$ of $U$, define
\[
 c_\ell(s)=\min\left\{\max(N,\deg_L Z^+,\deg_L Z^-):
 \begin{array}{l}N\in\Z_{>0},\quad Z^\pm\text{ effective integral cycles},\\
 \cl_\ell(Z^+-Z^-)=Nt_{\ell,s}\end{array}\right\},
\]
with value $\infty$ if the set is empty. Write $c_\eta$ for the analogous
geometric characteristic-zero minimum in rational Betti cohomology.
Comparison makes equality of these rational generic classes in one
$\ell$-adic realization equivalent to equality in Betti cohomology. If the
local condition requires one pair of cycles to give equality in all
compatible realizations, write $c_{\rm all}(s)$; then
$c_\ell(s)\leq c_{\rm all}(s)$.

\subsection{A proper parameter space with a closed class condition}

\begin{proposition}[B01: bounded effective pairs]
For each integer $D\geq0$ there is a projective finite-type $U$-scheme $H_D$
whose geometric points parameterize pairs of effective codimension-$r$
cycles of $L$-degrees at most $D$, and every such pair occurs. For each
$N\geq1$, the class condition
\[
 \cl_\ell(Z^+-Z^-)=Nt_\ell
\]
defines an open-and-closed subscheme $H_{N,D}(t_\ell)$ of $H_D$. In particular,
this subscheme is projective over $U$.
\end{proposition}

\begin{proof}
The Hilbert polynomials of reduced integral $d$-dimensional subschemes of
geometric fibers having $L$-degree at most $D$ form a finite set
$\mathcal P_D$. This is the boundedness assertion for reduced equidimensional
subschemes of bounded degree in Grothendieck's FGA, Bourbaki Expos\'e~221,
Lemma~2.4. Each $\operatorname{Hilb}^{P}(X/U)$ is projective over $U$ by
Theorem~3.2 of the same expos\'e. If $P$ is computed with $L^m$, let
$\delta(P)$ be $d!$ times its leading coefficient divided by $m^d$.

An effective cycle of degree at most $D$ has the form
$\sum_{i=1}^q n_i[Y_i]$, with $q\leq D$, $1\leq n_i\leq D$, and
$\sum_i n_i\deg_L(Y_i)\leq D$. Let $Q_D$ be the finite disjoint union of
products
\[
 \prod_{i=1}^q\operatorname{Hilb}^{P_i}(X/U)
 \quad\text{with}\quad
 \sum_{i=1}^q n_i\delta(P_i)\leq D,
\]
over all such finite lists of polynomials and multiplicities. Include the
empty product $U$ for the zero cycle. Set $H_D=Q_D\times_U Q_D$. A finite
disjoint union and a finite product of projective $U$-schemes are projective.

The universal subschemes on the Hilbert factors are flat over their parameter
schemes. At an arbitrary boundary point their fibers need not be reduced or
pure, but their Hilbert polynomial still has top degree $d$. Their
top-dimensional fundamental cycles are effective and have degree
$\delta(P_i)$. Lower-dimensional and embedded components do not affect that
cycle or its leading Hilbert coefficient. Thus every point of $H_D$ gives a
valid pair of cycles with the claimed degree bounds, and every desired pair
occurs by choosing its reduced integral supports.

The relative cycle-class construction supplies sections of the pullback of
$\mathcal V_\ell$ for these flat universal cycles. The precise source is
SGA~$4\frac12$, \emph{La classe de cohomologie associ\'ee \`a un cycle},
\S\S2.3.1--2.3.3, Lemma~2.3.7 and Theorem~2.3.8. The hypothesis is fiber
dimension at most $d$, not purity of every fiber. The more general
construction uses a weight of finite Tor dimension over the base; the
structure sheaf of a flat universal subscheme qualifies. Smoothness is
required of the ambient $X$ over the parameter scheme, which is preserved
by base change, rather than of the parameter scheme itself. Arbitrary
base-change compatibility identifies the stalk with the ordinary cycle
class of the fiber's top-dimensional fundamental cycle. Consequently
singular and nonflat parameter schemes over $U$ cause no difficulty.

Subtract the negative universal class and $N$ times the pullback of
$t_\ell$. The resulting global section $\Delta$ of a lisse $\Q_\ell$-sheaf
vanishes at one geometric point of a connected component if and only if it
vanishes on that component. This follows from the fundamental-group
description of a lisse sheaf; equivalently, use an integral lattice and
reduce modulo every $\ell^a$. A noetherian scheme has finitely many connected
components, each open and closed. The zero condition is therefore an
open-and-closed union of these components, giving $H_{N,D}(t_\ell)$.
\end{proof}

\begin{remark}
Smooth proper base change by itself does not construct the required class
locus on an unspecified Chow parameter space. The finite Hilbert
parameterization and the relative fundamental class supply that missing
step. The base-change reference is SGA~4, Expos\'e~XVI, Corollary~2.2; the
locally constant coefficient hypothesis omitted in the original printing
is imposed here and is automatic for the constant coefficients used above.
\end{remark}

\subsection{The precise fixed-multiplier criterion}

\begin{theorem}[B02: corrected form of sketch Lemma~2.5]
Fix $N\geq1$ and $D\geq0$. The following conditions are equivalent:
\begin{enumerate}
\item On $X_{\overline K}$ there are effective integral cycles $Z^\pm$ of
$L$-degrees at most $D$ with $\cl_B(Z^+-Z^-)=Nt$.
\item Such cycles with $\ell$-adic difference class $Nt_{\ell,v}$ exist on
$X_{\overline{k(v)}}$ for infinitely many $v\in U$.
\item Such cycles exist for every finite place $v\in U$.
\end{enumerate}
The equivalence also holds when the local condition requires the same pair
of cycles to realize all compatible specializations of $t$.
\end{theorem}

\begin{proof}
The condition at a point is precisely nonemptiness of the corresponding
geometric fiber of $H_{N,D}(t_\ell)\to U$. Its image is closed by properness.
A proper closed subset of the noetherian integral one-dimensional scheme
$U$ is finite. Thus infinitely many represented closed fibers force the
generic point into the image. A nonempty finite-type generic fiber has a
closed point over a finite extension $K'/K$. Its cycles retain the same
integer multiplicities and degree bounds, and their $\ell$-adic difference
is $Nt_\ell$. Embed $K'$ into $\C$ extending the original embedding.
Betti--\'etale comparison and the injection of a rational vector space into
its scalar extension to $\Q_\ell$ give the Betti equality in (1).

Conversely, a generic pair in (1) gives a point of the generic class locus.
The closed image of that locus contains the generic point of $U$, hence is
all of $U$, proving (3). The implication (3)$\Rightarrow$(2) is immediate.

For the simultaneous-realization statement, the backward implication needs
only the one fixed $\ell$ already used. In the other direction, define the
actual characteristic-zero cycles over a finite extension and specialize
their Hilbert points along each good trait, using properness. The resulting
top-dimensional cycles remain effective, with the same degrees and
multiplicities. Compatibility of cycle classes with specialization gives
the target equality in every chosen realization. This constructs one pair
of cycles, rather than unrelated choices for different $\ell$.
\end{proof}

\begin{remark}[B02a: rational versus integral algebraicity]
For prescribed $N$, the assertion that ``$Nt$ is algebraic'' must mean that
$Nt$ is the class of an integral cycle. If it means membership in the image
of $\CH^r\otimes\Q$, the literal fixed-$N$ statement is false. On
$\Pp^g$, take $t=h^r/2$ and $N=1$. This class is rationally algebraic, but
the class of every integral cycle difference is an integer multiple of
$h^r$, in every characteristic. Either retain the integral interpretation
or quantify existentially over $N$.

The general bounded-rational-complexity principle has precedents:
Br\"unjes--Serp\'e, \emph{Uniform bounds and ultraproducts of cycles},
Definition~2.3(ii) and Theorem~3.4, prove such a criterion at infinitely
many closed points. The number-field extension is discussed immediately
before their \S2. The argument here provides the fixed-$(N,D)$ geometric
fiber form and the proper parameter space needed for subsequent claims;
it does not claim the basic principle as new.
\end{remark}

\subsection{Exact stabilization of the least complexity}

\begin{theorem}[B03: eventual constancy or cofinite divergence]
If $c_\eta=C_0<\infty$, then $c_\ell(v)\leq C_0$ for all $v\in U$, and
\[
 c_\ell(v)=c_{\rm all}(v)=C_0
\]
outside a finite subset of $U$, with the simultaneous assertion under the
compatibility assumptions above. If $c_\eta=\infty$, then for every finite
$C$ the set
\[
 \{v\in U:c_\ell(v)\leq C\}
\]
is finite. Thus both local minima tend to infinity in the cofinite filter.
\end{theorem}

\begin{proof}
For an integer $C\geq1$, the finite disjoint union
\[
 \coprod_{1\leq N\leq C}H_{N,C}(t_\ell)
\]
is projective over $U$. Its image $F_C$ is closed, and its generic fiber is
nonempty exactly when $c_\eta\leq C$. If $c_\eta=C_0$, then $F_{C_0}=U$,
whereas $F_{C_0-1}$ is finite when $C_0>1$. Outside that finite set no smaller
value is possible. When $C_0=1$ there is no smaller allowed complexity.

A generic triple attaining $C_0$ specializes as an actual pair of cycles
and gives $c_{\rm all}(v)\leq C_0$. The inequalities
$c_\ell(v)\leq c_{\rm all}(v)$ supply the simultaneous assertion. Finally,
if $c_\eta=\infty$, every $F_C$ is a proper closed subset of $U$, hence finite.
\end{proof}

\begin{remark}
For a fixed polarized arithmetic abelian variety and a fixed algebraic
class, the eventual constant is the same along the ordinary and the
supersingular subsequences whenever those subsequences are infinite.
Degrees or denominators attached to a chosen CM construction may grow
while the minimum remains constant. The theorem determines this shape of
the function; by itself it computes neither $C_0$ nor the exceptional
primes. In the alternative $c_\eta=\infty$, local algebraicity at every
prime, if separately known, means finite values escaping every fixed bound.
\end{remark}

\subsection{A class-independent exceptional integer}

\begin{proposition}[B04: one-prime criterion]
For each $D$ there is a nonzero integer $b_D$, depending only on the
polarized model, codimension, and the chosen finite parameterization
$H_D$, such that the following holds. If $v\nmid b_D$ and effective cycles
of degrees at most $D$ represent $Nt_{\ell,v}$ at $v$, then effective
cycles with the same degree bounds represent $Nt$ over $\overline K$.
No bound on $N$ is required. In particular, if $c_\eta>C$, every place
with $c_\ell(v)\leq C$ divides $b_C$.
\end{proposition}

\begin{proof}
Generic flatness, EGA~IV$_2$, Theorem~6.9.1, gives a nonempty open subset
over which $H_D\to U$ is flat. Remove its finite complement by inverting
one nonzero rational integer $b_D$. A point $z\in H_{N,D}(t_\ell)$ over a
remaining place has a generization $z_\eta$ lying over the generic point:
flat morphisms satisfy going down, applied to a flat
$\mathcal O_{K,v}$-algebra in affine neighborhoods. Since the class locus
is open and closed in $H_D$, it contains every generization of $z$.
Its generic fiber is therefore nonempty. Choose a closed point of that
fiber and use the comparison argument of B02. The integer $N$ and degree
bounds have not changed.
\end{proof}

\paragraph{Making $b_D$ concrete.}
Suppose equations and the finite Hilbert-polynomial list defining $H_D$
are available. Cover it by finitely many affine schemes
$\Spec A_i$ over $R=\mathcal O_K[S^{-1}]$. The $R$-torsion ideal
\[
 T_i=\ker(A_i\longrightarrow A_i\otimes_RK)
\]
is finitely generated as an $A_i$-ideal. Choose a nonzero element $a_i\in R$
annihilating all of its finitely many generators. After inverting
$\prod_i a_i$, every $A_i$ is torsion-free over the Dedekind domain $R$,
and is therefore flat. Clearing denominators and using a rational integer
divisible by the norm gives a possible $b_D$. This procedure uses the full
parameter scheme and requires no decision procedure for equality to $t$.
Effective ideal arithmetic can implement it when those explicit finite
presentations are supplied. No numerical $b_D$ or general size estimate is
computed here.

For a construction with degree bound $D(v)$, one prime satisfying
$v\nmid b_{D(v)}$ suffices. Infinitely many primes with a fixed $D$ ensure
such a prime. When $D(v)$ varies, a prime-counting theorem alone does not
compare the selected prime to the exceptional integer indexed by its own
degree bound.

\subsection{Descent and automatic multiplier boundedness}

\begin{proposition}[B05: descent without increasing complexity]
Let $X,L$ be defined over an algebraically closed subfield $k\subset\C$.
A pair of complex effective cycles of degrees at most $D$ whose rational
Betti difference is $Nt$ can be replaced by a pair defined over $k$ with
the same degree bounds and Betti class. In particular, for $k=\overline K$
the complex and geometric generic minima agree.
\end{proposition}

\begin{proof}
Use the finite product of Hilbert schemes corresponding to the supports
and multiplicities of the given pair. A $\C$-point lies on a connected
component. A connected component of a projective finite-type scheme over
the algebraically closed field $k$ is geometrically connected, so the
chosen component is the base change of a component over $k$. It has a
$k$-point. On this parameter space the universal class is a section of
the constant cohomology sheaf of $X$, hence the two points give the same
$\ell$-adic class. Degree and multiplicities are unchanged. Comparison
gives the rational Betti equality.
\end{proof}

\begin{proposition}[B06: bounded numerator degrees already bound multipliers]
Suppose $t_\ell\ne0$. For each fixed $D$ there is a finite set
$M_D\subset\Z_{>0}$ containing every multiplier $N$ that can occur in
\[
 \cl_\ell(Z^+-Z^-)=Nt_{\ell,s},\qquad \deg_L Z^\pm\leq D,
\]
at any geometric fiber $s$ of $U$. Consequently rational algebraicity of
$t$ is equivalent to bounded degrees of the two integral numerators at
infinitely many good primes; no separate multiplier bound is needed.

The analogous assertion holds for a smooth projective family over a
connected noetherian complex base carrying a nonzero flat rational class
$t$: bounded integral numerator degrees on a Zariski-dense set force a
uniform complexity bound and algebraicity at every fiber.
\end{proposition}

\begin{proof}
Let $C_1,\ldots,C_q$ be the connected components of $H_D$. Write $\alpha$
for the universal difference class and $\beta$ for the pullback of
$t_\ell$. Every stalk of $\beta$ is nonzero, since $U$ is connected and
$t_\ell$ is a nonzero global lisse section. On any one $C_i$, at most one
integer $N$ can satisfy $\alpha=N\beta$: two would give
$(N-N')\beta=0$ in a $\Q_\ell$-vector space. Equality at one geometric
point already implies equality throughout $C_i$, by B01. The possible
integers therefore form a set $M_D$ of size at most $q$.

If degrees at most $D$ work at infinitely many primes, some $N\in M_D$
works infinitely often. Apply B02. Conversely a rational algebraic
class has the form $(\cl Z^+-\cl Z^-)/N$ for effective integral cycles,
and their specializations give a uniform numerator bound. For $t=0$
take both cycles zero and $N=1$; the algebraicity conclusion remains
trivial, although all positive multipliers are possible.

Over a connected noetherian complex base, repeat B01 with the Betti
local system, or its $\ell$-adic counterpart. The parameter scheme is
still noetherian with finitely many connected components. Each component
admits at most one multiplier. Their relevant union is a finite union of
open-and-closed components and is projective over the base. Its image is
closed. If it contains a Zariski-dense set, it is the whole base; the finite
list of multipliers supplies the uniform complexity bound.
\end{proof}

\begin{remark}
The numerator degrees in B06 count integer multiplicities. A bound on
reduced supports, on rationally normalized degrees, or on a denominator
alone is a different condition. B06 does not evaluate $\max M_D$; it
strictly reduces a boundedness argument to controlling the integral
numerator degrees.
\end{remark}

\subsection{Exact source locators and verification level}
The following locators concern the sources actually used in B01--B06.
They do not purport to certify the full foundations of FGA, EGA, or SGA.
\begin{itemize}
\item Grothendieck, FGA, \emph{Techniques de construction et th\'eor\`emes
d'existence en g\'eom\'etrie alg\'ebrique IV: les sch\'emas de Hilbert},
Bourbaki Expos\'e~221 (1961), Lemma~2.4, printed p.255; Theorems~3.1--3.2,
pp.259--260. The original statements were read at
\url{https://www.numdam.org/article/SB_1960-1961__6__249_0.pdf}.
\item Grothendieck--Dieudonn\'e, EGA~IV$_2$, \S6.9, Theorem~6.9.1 and
Lemma~6.9.2, Publ. Math. IHES~24 (1965), p.154. Canonical metadata:
\url{https://www.numdam.org/articles/10.1007/BF02684322/}.
The exact theorem was checked by indexed retrieval; the entire original
PDF exceeded the web retrieval size limit. Generic flatness is used as
background.
\item SGA~4, Expos\'e~XVI, Corollary~2.2, original p.214. Statement read in
the re-edition at \url{https://www.normalesup.org/~forgogozo/SGA4/16/16.pdf}.
The locally constant coefficient hypothesis from the SGA~$4\frac12$
erratum, p.312, is explicitly imposed.
\item Grothendieck, r\'edig\'e par Deligne, \emph{La classe de cohomologie
associ\'ee \`a un cycle}, SGA~$4\frac12$, LNM~569 (1977), pp.129--153,
\S\S2.3.1--2.3.3, Lemma~2.3.7, Theorem~2.3.8 and Remark~2.3.10.
Chapter metadata: \url{https://link.springer.com/chapter/10.1007/BFb0091521}.
The relative construction and base-change statements were read in the
retypeset source linked by Milne,
\url{https://github.com/NomiL/sga4.5/blob/master/sga4.5-4.tex}, and checked
against indexed text of
\url{https://math.bu.edu/people/yangzhe/Cohomologie_Etale.pdf}, p.70.
No claim of re-proving the full cycle-class theory is made.
\item Deligne, notes by Milne, \emph{Hodge cycles on abelian varieties},
LNM~900 (1982), Proposition~2.9(a)--(b), revised author-hosted text,
pp.18--19: \url{https://www.jmilne.org/math/Documents/Deligne82.pdf}.
The finite-Galois-orbit statement and its three proof alternatives were
read.
\item Br\"unjes--Serp\'e, \emph{Uniform bounds and ultraproducts of cycles},
arXiv:0901.4853, Definition~2.3(ii), Theorem~3.4, and the introduction's
number-field discussion: \url{https://arxiv.org/pdf/0901.4853}. The theorem
and proof were read; its nonstandard foundational dependencies remain
background.
\end{itemize}
